\chapter{Privilege, traps and events}\label{ch:traps}

\IS has two privilege modes, a small set of banked trap registers, and no interrupts. A \term{trap}
is the synchronous transfer of control that occurs when an instruction cannot complete. The only
asynchronous event that changes a core's control flow is a \term{stop request} sent through the
core-control space, and it is delivered as a trap at an instruction boundary. Everything else a core
waits for, including device events, arrives as the completion of a transfer or as a capability in its
inbox.

\section{Privilege modes}\label{sec:modes}

Each core runs in \term{supervisor mode} or \term{user mode}. Cores leave reset in supervisor mode.
The current mode and the mode that \insn{eret} returns to are recorded in CSR \csr{status}:

\begin{center}\small
\begin{tabularx}{\linewidth}{@{}c l X@{}}
\toprule
\textbf{Bit} & \textbf{Name} & \textbf{Meaning}\\\midrule
0 & S & 1 while the core runs in supervisor mode. Read-only: it changes only on trap entry and \insn{eret}.\\
1 & SPP & The mode a supervisor \insn{eret} enters: 1 supervisor, 0 user. Set by trap entry; writable by the supervisor.\\
63:2 & --- & Read as zero; writes are ignored.\\
\bottomrule
\end{tabularx}
\end{center}

A core enters user mode only by executing \insn{eret} in supervisor mode with SPP = 0. It returns to
supervisor mode only by taking a trap that is routed to the supervisor. User mode differs from
supervisor mode in exactly the operations of Table~\ref{tab:userlimits}; every other instruction,
including every transfer, far atomic, overlay operation and the endpoint operations \insn{csend} and
\insn{crecv}, behaves identically in both modes. Authority over memory and devices is never a matter of mode: it comes from the capabilities
a core holds.

\begin{table}[htbp]
\centering\small
\caption{What user mode cannot do}\label{tab:userlimits}
\begin{tabularx}{\linewidth}{@{}l X@{}}
\toprule
\textbf{Operation} & \textbf{In user mode}\\\midrule
\insn{csave}, \insn{crestore}, \insn{csavex}, \insn{crestorex} & \cause{priv}: the capability save table belongs to the supervisor\\
\insn{cdom}, \insn{crevoke}, \insn{cundom}, \insn{cmkep} & \cause{priv}: revocation and endpoints are the supervisor's\\
\insn{csrr} of \csr{sscratch} or \csr{ipend} & \cause{priv}\\
\insn{csrw} of \csr{status}, \csr{corebase}, \csr{corecount}, \csr{ucause}, \csr{uepc}, \csr{utval}, \csr{utvec}, \csr{sscratch}, \csr{ipend} & \cause{priv}\\
\insn{ctake}, \insn{cpoll}, \insn{tany} with bit~16 set & \cause{priv}: the inbox belongs to the supervisor (Section~\ref{sec:userinbox})\\
\insn{halt} & after draining the core's transfers, traps to the supervisor with cause \cause{halt} instead of stopping the core\\
\insn{svc} & traps to the supervisor with cause \cause{svc} (in supervisor mode it is \cause{illegal})\\
\insn{cgive} & restricted to the core window: \cause{capop} if \csr{corecount} is 0\\
\bottomrule
\end{tabularx}
\end{table}

\subsection{Banked trap state}

\csr{cause}, \csr{epc}, \csr{tval} and \csr{tvec} exist twice: once for the supervisor and once for
user mode. An access by name reaches the copy of the current mode, so the same handler code works in
either mode. The supervisor reaches the user copies through \csr{ucause}, \csr{uepc}, \csr{utval} and
\csr{utvec} (CSRs 19--22), for example to install or inspect a program's handler. \csr{sscratch} is a
scratch register only the supervisor can read or write; \csr{scratch} is shared by both modes.

\begin{center}\small
\begin{tabular}{@{}l l l@{}}
\toprule
\textbf{Register} & \textbf{In supervisor mode} & \textbf{In user mode}\\\midrule
\csr{cause}, \csr{epc}, \csr{tval}, \csr{tvec} & supervisor copy & user copy\\
\csr{ucause}, \csr{uepc}, \csr{utval}, \csr{utvec} & user copy & user copy, read-only\\
\csr{sscratch}, \csr{ipend} & read and write & \cause{priv}\\
\csr{scratch} & read and write & read and write\\
\bottomrule
\end{tabular}
\end{center}

The pseudocode writes the two copies as \texttt{SEPC}, \texttt{STVEC}, \ldots\ and \texttt{UEPC},
\texttt{UTVEC}, \ldots

\section{Trap entry and routing}\label{sec:trapentry}

A trap in supervisor mode always goes to the supervisor handler. A trap in user mode goes to the
program's own handler when the program has installed one and the cause is the program's business;
otherwise it goes to the supervisor:

\begin{pseudo}
func Trap(c, value)
    if Mode == User and c not in {halt, stop, priv, svc}
       and UTVEC != 0 and PC != UTVEC then
        UEPC = PC; UCAUSE = c; UTVAL = value;     // the program's own handler,
        NextPC = UTVEC;                           // still in user mode
    else
        double = (Mode == Supervisor and STVEC != 0 and PC == STVEC);
        SEPC = PC; SCAUSE = c; STVAL = value;
        STATUS.SPP = if Mode == User then 0 else 1;
        Mode = Supervisor;
        if STVEC == 0 or double then
            Halted = TRUE;                        // reported as a (double) fault
        else
            NextPC = STVEC;
        end
    end
end
\end{pseudo}

\begin{itemize}
\item \cause{halt}, \cause{stop}, \cause{priv} and \cause{svc} always reach the supervisor: a program
  can neither catch its own end nor its own preemption, nor hide a privilege violation, nor answer its
  own supervisor call.
\item A fault on the first instruction of the user handler ($\reg{pc} = \csr{utvec}$) goes to the
  supervisor, so a broken user handler cannot loop.
\item In supervisor mode, a fault on the first instruction of the handler is a \term{double fault} and
  stops the core, as does any supervisor trap while \csr{tvec} is zero.
\end{itemize}

A trap is \term{precise}: the trapping instruction has had no architectural effect (no register, Local
Store, capability, far memory or device state has changed, and no transfer has started). It does not
retire: \csr{instret} is not incremented. Execution continues at the handler with all other state
unchanged.

\section{Causes}

\begin{table}[htbp]
\centering\small
\caption{Trap causes (summary; see Appendix~\ref{app:causes} for \csr{tval})}\label{tab:causes-short}
\begin{tabularx}{\linewidth}{@{}r >{\sffamily}l X@{}}
\toprule
\textbf{Code} & \textrm{\bfseries Cause} & \textbf{Meaning}\\\midrule
1 & illegal & undefined encoding\\
2 & lsfault & Local Store address out of range or misaligned\\
3 & capinv & capability register or save slot is not live (empty, or revoked)\\
4 & perm & capability lacks the permission or names the wrong kind of space\\
5 & bounds & far range outside the segment\\
6 & xalign & memory transfer or atomic not 8-byte aligned; \insn{xmap} of a unit not aligned to its slot size\\
7 & ifault & instruction fetch outside the I-LS or misaligned\\
8 & xmiss & \insn{xcall} target, or \insn{xret} far link, not in the overlay map\\
9 & pend & access to an in-flight range with CCR.PTRAP set (debug)\\
10 & capop & capability operation precondition failed (overwrite, split, join, give, revocation)\\
11 & priv & privileged operation in user mode\\
12 & halt & \insn{halt} in user mode\\
13 & stop & a stop request sent through core control, taken in user mode\\
14 & svc & \insn{svc} in user mode: a supervisor call\\
\bottomrule
\end{tabularx}
\end{table}

\section{Priority}

At most one trap is taken per instruction. The order of checks is:
\begin{enumerate}
\item a pending stop request, in user mode: \cause{stop};
\item instruction fetch: \cause{ifault};
\item decoding: \cause{illegal};
\item the instruction's own checks, in the order listed in its \emph{Exceptions} table in
  Part~\ref{part:isa} (for privileged instructions, \cause{priv} comes first).
\end{enumerate}
For example, a \insn{tget} through an empty capability with a misaligned length reports
\cause{capinv}, not \cause{xalign}.

\section{The trap value}

\csr{tval} identifies what was wrong:
\begin{itemize}
\item for transfer faults, the far offset \reg{ro}, except \cause{lsfault}, where it is the Local
  Store address;
\item for loads and stores, the effective address;
\item for \cause{ifault}, the faulting \reg{pc}; for \cause{illegal} and \cause{priv}, the instruction
  word;
\item for capability instructions, the register or slot number concerned, or the offending value (see each
  instruction);
\item for \cause{xmiss}, the absolute far address of the call target, or the far link (bit~63 set) of
  an \insn{xret};
\item for \cause{halt}, \cause{stop} and \cause{svc}, zero.
\end{itemize}

\section{Returning from a trap}

\insn{eret} continues at \csr{epc} of the current mode. In user mode it stays in user mode. In
supervisor mode it enters the mode that \csr{status}.SPP names, so a supervisor handler that took a
trap from a program returns to the program, and one that took a trap from supervisor code returns to
supervisor code. After a trap \csr{epc} holds the address of the trapping instruction: to resume after
it, add 4 to \csr{epc}; to retry it (as an overlay or paging handler does), return with \csr{epc}
unchanged. A supervisor starts a program the same way: it sets \csr{epc} to the entry point, clears
SPP and executes \insn{eret}.

\section{Supervisor calls}\label{sec:svc}

\insref{svc} is the one instruction whose purpose is to trap. Executed in user mode it enters the
supervisor with cause \cause{svc}, \csr{tval} = 0 and \csr{epc} = the address of the \insn{svc}; the
supervisor serves the call and returns with \csr{epc} + 4. It is never routed to the program's own
handler. In supervisor mode \insn{svc} is an illegal instruction: a supervisor has nobody to call.

The architecture defines no arguments. A supervisor reads the program's registers, which it has just
saved, and writes results back into the saved copies. Supervisor calls are few in \IS because
authority comes from capabilities, not from the supervisor: a program reads a file by sending a message
to a server through an endpoint, which is a user-mode instruction. What only the supervisor of the
program's core can do is \emph{stop running the program until something happens}, and give it what
arrived in the core's inbox meanwhile. StrataOS has eleven supervisor calls, of which blocking for a
reply, receiving a request and yielding are the important ones (Section~\ref{sec:os-svc}).

\begin{asm}
        csend c5, c2                    ; lend the message buffer to a server (user mode)
        li    a0, 1                     ; StrataOS convention: a0 = the function
        svc                             ; "kernel: wake me when it comes back"
        bltz  a0, failed                ; the supervisor wrote the result into a0
\end{asm}

\section{Writing a handler}

A handler starts with no free register. A supervisor handler frees one with \csr{sscratch}, which user
code can neither read nor write; this is how the StrataOS handler saves a program's registers without
trusting anything the program left behind:

\begin{asm}
w_trap:                                 ; the supervisor's tvec
        csrw  sscratch, t0              ; free t0 (only the supervisor can see sscratch)
        la    t0, w_save                ; t0 = register save area in kernel data
        sd    x1, 8(t0)                 ; save x1 ... x31
        ...
        csrr  t1, sscratch              ; the program's t0
        sd    t1, 40(t0)                ; save it too
        la    sp, kstack_top            ; a stack of our own
        csrr  t1, cause                 ; the supervisor's copy: why are we here?
\end{asm}

Bare-metal code and user handlers save registers to a fixed D-LS area addressed from \reg{x0} (offsets
up to 32767), or stash one register in \csr{scratch}. A handler that may itself trap must save
\csr{epc}, \csr{cause} and \csr{tval} first, because a nested trap in the same mode overwrites them.

\begin{asm}
handler:                                ; tvec points here
        sd    t0, hsave+0(zero)         ; save t0 without needing a base register
        csrr  t0, cause
        sd    t0, last_cause(zero)      ; remember why
        csrr  t0, epc
        addi  t0, t0, 4                 ; skip the trapping instruction
        csrw  epc, t0
        ld    t0, hsave+0(zero)
        eret
\end{asm}

The same handler works on the bare machine, where it runs in supervisor mode, and under StrataOS, where
it runs in user mode: there its \insn{csrw tvec} installs the user handler and the banked registers it
reads are the user copies.

\section{Halting}\label{sec:halt}

\insn{halt} first waits for all of the core's in-flight transfers to complete, pending device reads
included, so output a program issued before halting is not lost. The drain happens in both modes: in
user mode it precedes the \cause{halt} trap. A core must therefore not halt while it has a device read
in flight that may never complete; a read of the console break port is retired first by writing the
release port (Section~\ref{sec:break}). In supervisor mode the core then stops for good; \reg{a0} is
conventionally its exit status and simulators report it. A halted core's capabilities are destroyed,
and any later \insn{cgive} to it traps with \cause{capop}.

In user mode \insn{halt} does not stop the core: it traps to the supervisor with cause \cause{halt}
and \csr{epc} = the address of the \insn{halt}. A program written for the bare machine therefore runs
unmodified under an operating system, which regains control when it ends.

A core also stops when a trap reaches the supervisor while \csr{tvec} is zero, or on a double fault.
The machine stops when every core has stopped. A machine in which every live core is waiting in
\insn{cgive}, \insn{ctake}, \insn{csend}, \insn{crecv}, \insn{tany} or \insn{halt}, and in which no
pending device read can still complete, is deadlocked; the simulators detect and report this.

\section{Stopping a core: core control}\label{sec:corectl}

Core-control space (space 5) is the only way one core can interrupt another. It has one doubleword per
core: offset $8n$ belongs to core $n$. A \insn{tput} through a core-control capability that writes a
doubleword to offset $8n$ acts on core $n$ according to its low two bits: bit~0 sends a \term{stop
request}, bit~1 sends a \term{reset request} (below); a doubleword with both clear has no effect. At
reset core~0 holds \reg{c4}, a copyable write-only capability for the whole space, $[0, 8N)$.
Core-control capabilities permit only \insn{tput}: \insn{tget} and \insn{tgetx} through them trap with
\cause{perm}.

A \term{reset request} restarts the target core as if it had just been powered on, but leaves the
shared memories and devices alone. Unlike a stop request it is taken \emph{in any mode}, also by a
core that has halted, so a core whose supervisor will not yield can still be stopped. A reset clears
the core's integer and capability registers (\reg{sp} gets \texttt{DLS\_SIZE} again, as at power-on),
its capability save table, its inbox, its overlay map, its transfer tags, and its CSRs (\csr{instret}
counts from the reset); reloads the \term{boot ROM}, the first 4~KiB of each Local Store (all of a
4~KiB store), from the boot image and clears the rest of both Local Stores; and
restarts the core at the boot ROM entry (\reg{pc} 0) in supervisor mode with the reset capabilities of
Table~\ref{tab:resetcaps}. It then runs the boot ROM's secondary path and waits for the system core to
hand it the operating system again, exactly as a core does at power-on. Whoever holds core control
over a core may reset it: the same right that lets them stop it. An operating system uses this to
reclaim a core from a guest that runs its own supervisor (Section~\ref{sec:os-guests}).

Precisely: the target's inbox is emptied when the request arrives, so a gift made after the
\insn{tput} that carried it (a \insn{cgive} waits for that \insn{tput}, Section~\ref{sec:publication})
reaches the restarted core. The target itself takes the request at its next instruction boundary, or
while an instruction waits, without doing anything of that instruction: every transfer it issued
before has then taken effect (they complete at issue), and a device read that is still waiting for its
device (console input, the break port, WAIT, VSYNC, EVENT, RX) is abandoned, as if never issued. A
core that writes a reset request to itself takes it before its next instruction. In hardware the
restart is not instantaneous: the core rewrites its Local Stores one doubleword per cycle (as many
cycles as the larger store has doublewords: 8192 with 64~KiB) before its first fetch; gifts sent to it
meanwhile wait in its inbox. The RTL keeps
the boot ROM in a small read-only memory beside each Local Store (that is why only 4~KiB come back:
a program image larger than that loses the rest at a reset, so the code a reset core runs must lie in
the first 4~KiB of each store, as the boot ROM does).

More generally, a \insn{tput} through a core-control capability takes effect before the core that
issued it starts its next instruction, so a stop or reset request a core sends is never overtaken by
what it does next.

A stop request sets the target core's \csr{ipend} bit. A core in user mode takes it before its next
instruction, and also while an instruction is waiting (in \insn{ctake}, \insn{crecv}, \insn{cgive},
\insn{csend}, \insn{tany}, \insn{twait} or a stall): \csr{ipend} is cleared and the core traps to the
supervisor with cause \cause{stop}, \csr{tval} = 0 and \csr{epc} = the instruction that did not
execute. A waiting instruction is abandoned with no effect, so \insn{eret} would simply resume the
wait. A core in supervisor mode is never interrupted: the request stays pending in \csr{ipend} until
the core next runs in user mode. The supervisor may read \csr{ipend}, clear it to discard a request, or
set it.

This is the only forced preemption in the architecture. There is no timer interrupt: the timer is a
device (Section~\ref{sec:timer}), and a core learns that a deadline has passed because its read of the
WAIT register completes. An operating system preempts by dedicating a core to this: the StrataOS system
core waits for the timer and, every time slice, writes a stop request to each worker core that has
more than one runnable process (Section~\ref{sec:os-preempt}). The same mechanism serves as a doorbell:
a core that has sent a message stops the destination core so that its supervisor looks at its inbox at
once.

\section{Events without interrupts}\label{sec:events}

\IS has no interrupt lines, interrupt controller, interrupt masks or interrupt handlers. A device
``interrupts'' by \emph{completing a transfer}: the core asks for the event by issuing a transfer that
the device completes only when the event happens, and learns of it the way it learns of any transfer
completing. Device-to-core notifications that carry authority arrive the same way capabilities always
do, as gifts in the inbox.

\begin{center}\small
\begin{tabularx}{\linewidth}{@{}l X@{}}
\toprule
\textbf{Event} & \textbf{How a core waits for it}\\\midrule
console input byte & a \insn{tget} from the console input port completes when the byte arrives\\
Ctrl-C & a \insn{tget} from the console break port completes when Ctrl-C is pressed (Section~\ref{sec:break})\\
a deadline & a \insn{tget} from the timer's WAIT register completes when NOW $\ge$ DEADLINE (Section~\ref{sec:timer})\\
the next display refresh & a \insn{tget} from the framebuffer controller's VSYNC register completes at the refresh (Section~\ref{sec:fbctl})\\
block-device command done & the buffer capability comes back to the inbox (Section~\ref{sec:blkdev})\\
message from another core & a capability arrives in the inbox\\
\bottomrule
\end{tabularx}
\end{center}

\insref{tany} \texttt{rd, rs} waits for the first of several such events. Bits 15:0 of \reg{rs}
select transfer tags, bit~16 selects the inbox; \insn{tany} returns as soon as a selected tag is idle
or a gift is waiting, with the mask of ready sources in \reg{rd}. The inbox bit is available in
supervisor mode only. The StrataOS system core sleeps on a pending timer read and its inbox at once:

\begin{asm}
        tget  tg11, t0, c13, t1, t2     ; read the timer's WAIT: completes at the DEADLINE
        ...
ev_wait:
        li    t0, (1 << 16) | (1 << 11) ; a gift, or tag 11 idle
        tany  t1, t0                    ; sleep until something happens
        srli  t2, t1, 16                ; a gift?
        bnez  t2, ev_msg                ; yes: take and serve it
        ...                             ; no: tg11 completed, the time slice is over
\end{asm}

Because an event is a transfer, everything the architecture says about transfers applies to it. The
read is checked and authorised by a capability when it issues; while it is pending its tag is busy and
loads from its D-LS bytes, and transfers that overlap them, wait, so a program that reads the event
byte early simply waits for it; and
nothing asynchronous ever runs on the core, so there are no handler races, no reentrancy and no
interrupt-masking discipline. Whether an event has happened \emph{yet} is timing-visible (the result of
\insn{tany} or \insn{tpoll}), just as for any transfer.

\subsection{The console break port}\label{sec:break}

The console's break port (device offset \hex{0010}) delivers Ctrl-C. A \insn{tget} from it stays in
flight until Ctrl-C is pressed and then completes with the byte 3. Ctrl-C presses that arrive while no
break read is pending are remembered and complete later reads at once, one press per read. A Ctrl-C typed at
the console never appears on the console input port.

\begin{rule*}{One release completes every pending break read}
A write to the release port (device offset \hex{0018}) completes, with zeros, every break read that is
waiting at that moment, on every core, and every read of the console input port that is waiting for its
first byte (all its bytes read zero, and no input is taken). A read issued after the release waits for
the next Ctrl-C, key or release. A release while no such read is waiting has no effect.
\end{rule*}

A program uses the release port to retire its break read when it no longer needs it, in particular
before it halts. Because one release retires the break reads of all cores, the cores that share the
console's break port must agree on who releases; a core whose break read completes with 0 instead of 3
knows it was a release. Software must keep at most one break read in flight per core. Both ports are
reached through the device root, so only a holder of a device capability covering them can wait for
Ctrl-C.
